\documentclass[11pt]{article}
\usepackage{amsmath, amssymb, geometry, mathtools, booktabs, longtable}
\usepackage{hyperref}
\geometry{margin=1in}

\title{RG Coarse-Lattice Fluctuation Flow for 2D $\phi^4$: Code-Level Description and CLI Reference}
\author{Kostas Orginos}
\date{\today}

\begin{document}
\maketitle

\section{Purpose}

This note documents the RG coarse-lattice fluctuation flow implemented in
\texttt{jaxqft/models/phi4\_rg\_coarse\_eta\_flow.py}, together with its
trainer
\texttt{scripts/phi4/train\_rg\_coarse\_eta\_flow.py}
and checker
\texttt{scripts/phi4/check\_rg\_coarse\_eta\_flow.py}.

This path is a new evolution branch. It leaves
\texttt{jaxqft/models/phi4\_mg.py},
\texttt{jaxqft/models/phi4\_rg\_cond\_flow.py},
\texttt{jaxqft/models/phi4\_rg\_checkerboard\_flow.py}, and
\texttt{jaxqft/models/phi4\_rg\_checkerboard\_inblock\_flow.py} unchanged.

\section{High-Level Design}

This architecture separates two conceptual steps:
\begin{enumerate}
\item RG blocking,
\item trivialization of the local fluctuations.
\end{enumerate}

At each non-terminal RG level:
\begin{itemize}
\item the current field is split into a coarse scalar field $c$ on the next
lattice and a local fluctuation vector $\eta\in\mathbb{R}^3$ on each coarse
site,
\item the coarse field $c$ is held fixed throughout the whole level update,
\item only $\eta$ is transformed,
\item the update of $\eta$ at one coarse site depends on the local coarse
field and neighboring coarse-lattice variables.
\end{itemize}

This is the key difference from the checkerboard branches:
\begin{quote}
there is no shifted reblocking and no attempt to preserve multiple incompatible
blockings. One blocking is chosen, its coarse field is fixed, and only the
fluctuations of that blocking are trivialized.
\end{quote}

\section{Lattice, Prior, and RG Split}

Let the fine lattice be
\[
\Lambda_0,\qquad |\Lambda_0| = L\times L,\qquad L=2^K.
\]
The batched field is
\[
\phi \in \mathbb{R}^{B\times L\times L}.
\]

The latent prior is sitewise standard normal:
\begin{equation}
p_Z(z)=\prod_{x\in\Lambda_0}\frac{1}{\sqrt{2\pi}}
\exp\!\left(-\frac{z_x^2}{2}\right).
\end{equation}

At RG level $k$, the current lattice has size
\[
L_k = L/2^k.
\]
Each $2\times 2$ block is packed into
\[
v=(v_{00},v_{01},v_{10},v_{11})\in\mathbb{R}^4.
\]

\subsection{Average Blocking}

For \texttt{rg\_type="average"}, the code uses the orthonormal basis
\begin{equation}
H=
\begin{pmatrix}
\tfrac12 & \tfrac12 & \tfrac12 & \tfrac12 \\
\tfrac12 & -\tfrac12 & \tfrac12 & -\tfrac12 \\
\tfrac12 & \tfrac12 & -\tfrac12 & -\tfrac12 \\
\tfrac12 & -\tfrac12 & -\tfrac12 & \tfrac12
\end{pmatrix}
\end{equation}
and defines
\begin{equation}
(c,\eta_1,\eta_2,\eta_3)=H\,v.
\end{equation}
Thus
\[
\eta=(\eta_1,\eta_2,\eta_3)
\]
spans the local fluctuation subspace.

The arithmetic block average is
\[
\bar v = \frac14\sum_{i\in\mathrm{block}} v_i = \frac12\,c.
\]
This arithmetic average is the coarse scalar used in conditioning.

\subsection{Select Blocking}

For \texttt{rg\_type="select"}, the code uses
\begin{equation}
c=v_{00},\qquad \eta=(v_{01},v_{10},v_{11}).
\end{equation}

\subsection{Inverse Reconstruction}

At each block the inverse reconstruction is
\[
v = R(c,\eta).
\]
For average blocking this is
\[
v = H^\top(c,\eta).
\]

\section{Recursive Structure}

The recursion stops at the final $2\times 2$ lattice, where the remaining
$4$-component block is mapped by an unconditional RealNVP.

For each non-terminal level:
\begin{enumerate}
\item split the current field into coarse field $c$ and local fluctuations
$\eta$,
\item recursively transform the coarse field,
\item transform $\eta$ conditional on the physical coarse field,
\item merge the transformed coarse field and transformed fluctuations.
\end{enumerate}

Thus if
\[
x_k \longleftrightarrow (c_{k+1},\eta_k),
\]
the level map is
\[
(z_{k+1},z_{\eta,k})
\mapsto
\bigl(x_{k+1},x_{\eta,k}\bigr)
\mapsto
x_k,
\]
where the fluctuation map is conditioned on $x_{k+1}$, not on the latent
coarse variable.

\section{Coarse-Lattice Fluctuation Field}

After blocking, one should think of the level-$k$ state as a field on the
coarse lattice
\[
\Lambda_{k+1},
\]
whose value at each site $b$ is
\[
y_b = (c_b,\eta_b),
\qquad
c_b\in\mathbb{R},\quad
\eta_b\in\mathbb{R}^3.
\]

The model does \emph{not} return to a fine lattice picture during the
fluctuation update. The fluctuation trivialization is performed directly on
this coarse lattice.

\section{Coarse-Lattice Neighborhood Conditioner}

Fix one RG level and suppress the level index. The conditioner input field is
\[
u_b = \bigl(\eta_b^{\mathrm{masked}}, \tilde c_b\bigr)\in\mathbb{R}^4,
\]
where
\[
\tilde c_b =
\begin{cases}
\tfrac12 c_b,& \text{average blocking,}\\
c_b,& \text{select blocking.}
\end{cases}
\]

The code uses a local patch MLP:
\begin{itemize}
\item choose a radius $r=\texttt{radius}$,
\item collect the square neighborhood
\[
\mathcal N_r(b)=\{b+\delta:\delta_y,\delta_x\in\{-r,\dots,r\}\},
\]
\item stack the $4$-component features from that neighborhood,
\item apply an MLP to the flattened patch.
\end{itemize}

Hence:
\begin{itemize}
\item \texttt{radius = 0} means on-site conditioning only,
\item \texttt{radius = 1} means a $3\times 3$ neighborhood on the coarse
lattice,
\item in general the receptive field contains $(2r+1)^2$ coarse-lattice sites.
\end{itemize}

\section{Red/Black Sweeps on the Coarse Lattice}

The coarse lattice is decomposed into red and black checkerboards:
\[
\Lambda_{k+1} = R \cup B,\qquad R\cap B = \varnothing.
\]
The corresponding masks are denoted by $M_R$ and $M_B$.

For one red sweep:
\begin{itemize}
\item all coarse variables $c$ are visible everywhere,
\item black fluctuations are visible and frozen,
\item red fluctuations are updated.
\end{itemize}

The masked fluctuation field for a red sweep is
\[
\eta^{\mathrm{masked}} = M_B \eta.
\]
So at a red site the conditioner sees:
\begin{itemize}
\item the local coarse scalar,
\item neighboring coarse scalars within the chosen radius,
\item neighboring black fluctuations within the chosen radius,
\item zero in the red fluctuation channels.
\end{itemize}

The black sweep is analogous, with red and black exchanged.

\subsection{Why This Remains Invertible}

The critical condition is:
\begin{quote}
when a site is updated, the conditioner must depend only on variables that are
held fixed during that substep.
\end{quote}

This is satisfied because:
\begin{itemize}
\item all coarse variables are fixed throughout the whole level map,
\item during a red sweep only black fluctuations are visible to the
conditioner,
\item during a black sweep only red fluctuations are visible to the
conditioner.
\end{itemize}

So increasing the radius does not break invertibility. It only enlarges the
neighborhood of frozen variables available to the conditioner.

\section{Per-Site Eta Map}

At each updated site $b$, the code applies a conditional lower-triangular
affine map on the local fluctuation vector
\[
\eta_b=(\eta_{b,1},\eta_{b,2},\eta_{b,3})\in\mathbb{R}^3.
\]

The conditioner outputs:
\begin{itemize}
\item three diagonal log-scales,
\item three strict lower-triangular entries,
\item three shifts.
\end{itemize}

Write these as
\[
\ell_{b,1},\ell_{b,2},\ell_{b,3},
\qquad
\lambda_{b,21},\lambda_{b,31},\lambda_{b,32},
\qquad
t_{b,1},t_{b,2},t_{b,3}.
\]
Then the sitewise forward map is
\begin{align}
x_{b,1} &= e^{\ell_{b,1}} \eta_{b,1} + t_{b,1}, \\
x_{b,2} &= \lambda_{b,21}\eta_{b,1} + e^{\ell_{b,2}}\eta_{b,2} + t_{b,2}, \\
x_{b,3} &= \lambda_{b,31}\eta_{b,1} + \lambda_{b,32}\eta_{b,2}
         + e^{\ell_{b,3}}\eta_{b,3} + t_{b,3}.
\end{align}

So the three components are mixed in one shot, but the Jacobian remains
triangular:
\[
\det J_b = e^{\ell_{b,1}+\ell_{b,2}+\ell_{b,3}}.
\]

\subsection{Inverse Map}

The inverse is solved exactly by forward substitution:
\begin{align}
\eta_{b,1} &= e^{-\ell_{b,1}}(x_{b,1}-t_{b,1}), \\
\eta_{b,2} &= e^{-\ell_{b,2}}
\bigl(x_{b,2}-t_{b,2}-\lambda_{b,21}\eta_{b,1}\bigr), \\
\eta_{b,3} &= e^{-\ell_{b,3}}
\bigl(x_{b,3}-t_{b,3}-\lambda_{b,31}\eta_{b,1}-\lambda_{b,32}\eta_{b,2}\bigr).
\end{align}

\subsection{Numerical Stabilization}

The code clips the conditioner outputs by
\[
\ell_{b,i}\leftarrow
\texttt{log\_scale\_clip}\cdot \tanh(\ell_{b,i}^{\mathrm{raw}})
\]
and
\[
\lambda_{b,ij}\leftarrow
\texttt{offdiag\_clip}\cdot \tanh(\lambda_{b,ij}^{\mathrm{raw}}).
\]

\section{$Z_2$ Symmetry}

The code supports
\[
\texttt{parity}\in\{\texttt{none},\texttt{sym}\}.
\]

For \texttt{parity="sym"}, the raw conditioner outputs are projected so that
the scale-like channels are even and the shifts are odd:
\begin{align}
\ell(u) &= \frac12\bigl(\hat \ell(u) + \hat \ell(-u)\bigr), \\
\lambda(u) &= \frac12\bigl(\hat \lambda(u) + \hat \lambda(-u)\bigr), \\
t(u) &= \frac12\bigl(\hat t(u) - \hat t(-u)\bigr).
\end{align}

Thus
\[
\ell(-u)=\ell(u),\qquad
\lambda(-u)=\lambda(u),\qquad
t(-u)=-t(u),
\]
which enforces the old-model-style
\[
z\mapsto -z,\qquad \phi\mapsto -\phi
\]
symmetry.

\section{One Full Level Map}

Let $P_R$ denote one red sweep and $P_B$ one black sweep on the coarse lattice.
One cycle is
\[
P = P_B \circ P_R.
\]
The code repeats this
\[
\texttt{n\_cycles}
\]
times per non-terminal level.

So the level fluctuation map is
\[
T_\eta = P^{N_c},
\qquad
N_c=\texttt{n\_cycles}.
\]

Unlike the checkerboard branches, there is no shifted reblocking here. The
coarse field of the chosen blocking remains fixed throughout the whole level
update.

\section{Terminal $2\times 2$ Flow}

At the coarsest $2\times 2$ lattice, the remaining four real variables are
packed into a vector in $\mathbb{R}^4$ and mapped by an unconditional RealNVP
\[
g_{\mathrm{term}}:\mathbb{R}^4\to\mathbb{R}^4.
\]

\section{Exact Log Density}

The inverse map returns
\[
(z,\log|\det J^{-1}|).
\]
Hence
\begin{equation}
\log p_\theta(x)
= \log p_Z(z)
+ \log|\det J_{\mathrm{term}}^{-1}|
+ \sum_k \log|\det J_k^{-1}|.
\end{equation}

At one updated site,
\[
\log|\det J_b^{-1}| = -(\ell_{b,1}+\ell_{b,2}+\ell_{b,3}).
\]
The log-Jacobian of one color sweep is the sum of this over the updated color.
The level log-Jacobian is the sum over all red and black sweeps in all cycles.

\section{Training Objective}

Training uses the standard flow objective
\begin{equation}
\mathcal L(\theta)=
\mathbb E_{z\sim p_Z}
\bigl[
\log p_\theta(g_\theta(z)) + S(g_\theta(z))
\bigr].
\end{equation}

For reweighting quality, the more relevant validation quantity is the centered
width of
\[
\Delta S = \log p_\theta(\phi) + S(\phi),
\]
namely
\[
\sigma\bigl(\Delta S-\langle\Delta S\rangle\bigr),
\]
because the reweighting factors are
\[
w=e^{-\Delta S}.
\]

\section{CLI Reference: Training Script}

This section documents
\texttt{scripts/phi4/train\_rg\_coarse\_eta\_flow.py}.

\begin{longtable}{p{0.24\linewidth}p{0.70\linewidth}}
\toprule
Argument & Meaning \\
\midrule
\endhead
\texttt{--resume} & Resume training from an existing checkpoint. \\
\texttt{--save} & Output checkpoint path. \\
\texttt{--save-every} & Save an intermediate checkpoint every $N$ epochs. Zero means save only at the end. \\
\texttt{--validate} & Run the final validation routine after training. \\
\texttt{--validate-batch} & Batch size used inside one validation draw. \\
\texttt{--validate-super} & Number of validation draws concatenated before computing $\Delta S$ and ESS summaries. \\
\texttt{--L} & Lattice linear size; the lattice is $L\times L$. Must be a power of two. \\
\texttt{--lam} & Quartic coupling $\lambda$. \\
\texttt{--mass} & Quadratic mass parameter $m^2$. \\
\texttt{--width} & Hidden width of the local patch MLP conditioner. \\
\texttt{--n-cycles} & Number of red/black sweep cycles per non-terminal RG level. \\
\texttt{--radius} & Coarse-lattice neighborhood radius used to build the local conditioner patch. \\
\texttt{--terminal-n-layers} & Number of RealNVP layers in the terminal $2\times 2$ flow. \\
\texttt{--terminal-width} & Hidden width of the terminal RealNVP MLPs. Zero means reuse \texttt{width}. \\
\texttt{--output-init-scale} & Small nonzero output initialization scale used to start near the identity while still allowing gradients to flow. \\
\texttt{--parity} & Symmetry mode. \texttt{sym} enforces the $\phi\to-\phi$, $z\to-z$ symmetry; \texttt{none} does not. \\
\texttt{--rg-type} & RG blocking rule: \texttt{average} or \texttt{select}. \\
\texttt{--log-scale-clip} & Tanh clip applied to the diagonal log-scales of the local triangular map. \\
\texttt{--offdiag-clip} & Tanh clip applied to the strict lower-triangular off-diagonal entries. \\
\texttt{--epochs} & Total target epoch count. On resume this is the final epoch number, not the number of extra epochs. \\
\texttt{--batch} & Training batch size. On resume, the checkpoint batch is reused unless this flag is explicitly passed. \\
\texttt{--lr} & Adam learning rate. On resume, the checkpoint learning rate is reused unless this flag is explicitly passed. \\
\texttt{--seed} & Random seed. \\
\bottomrule
\end{longtable}

\section{CLI Reference: Checker Script}

This section documents
\texttt{scripts/phi4/check\_rg\_coarse\_eta\_flow.py}.

\begin{longtable}{p{0.24\linewidth}p{0.70\linewidth}}
\toprule
Argument & Meaning \\
\midrule
\endhead
\texttt{--shape} & Lattice shape used for invertibility and timing tests. \\
\texttt{--jac-shape} & Small lattice shape used for the dense autodiff Jacobian check. \\
\texttt{--batch-size} & Batch size used in invertibility and timing tests. \\
\texttt{--width}, \texttt{--n-cycles}, \texttt{--radius}, \texttt{--terminal-n-layers}, \texttt{--terminal-width}, \texttt{--output-init-scale}, \texttt{--parity}, \texttt{--rg-type}, \texttt{--log-scale-clip}, \texttt{--offdiag-clip} & Same architectural meanings as in the training script. \\
\texttt{--lam}, \texttt{--mass} & Theory parameters used in the timing loss-gradient benchmark. \\
\texttt{--tests} & Comma-separated subset of \texttt{invertibility}, \texttt{jacobian}, \texttt{timing}. \\
\texttt{--invert-tol} & Absolute tolerance for the forward-then-inverse reconstruction test. \\
\texttt{--jac-tol} & Tolerance for the autodiff-vs-analytic logdet comparison. \\
\texttt{--jac-perturb-scale} & Small perturbation applied to initialized weights before the Jacobian test, so the logdet is not trivially zero. \\
\texttt{--timing-warmup} & Warmup iterations for each timed kernel. \\
\texttt{--timing-iters} & Number of timed iterations per kernel. \\
\texttt{--json-out} & Optional JSON output path for the checker summary. \\
\texttt{--selfcheck-fail} & Exit with nonzero status if any selected test fails. \\
\bottomrule
\end{longtable}

\section{Practical Interpretation}

This architecture is designed for the following principle:
\begin{quote}
blocking should define the coarse variables; the flow should then trivialize
only the fluctuations of that fixed blocking.
\end{quote}

So it gives up the shifted-blocking idea and instead invests model capacity in:
\begin{itemize}
\item preserving the physically chosen coarse field,
\item using neighboring coarse-lattice context,
\item mixing the three local fluctuation components exactly and cheaply.
\end{itemize}

Empirically this produces a strong quality-to-cost tradeoff:
\begin{itemize}
\item much cheaper than the transformer-heavy checkerboard branches,
\item cleaner mathematically than the shifted-blocking construction,
\item already competitive with the best previous models in reweighting quality.
\end{itemize}

\end{document}
